\answer{ex:17:1}{$0.2$~Вт проти $81$~мкВт; сирому потокові потрібно не більше $50$~пДж/біт.}
\answer{ex:17:2}{За $c=0.1$ $5.6$, $8.7$, $9.2$~біт/с, границя $9.3$~біт/с; за $c=0$ і $N=1000$ $65$~біт/с.}
\answer{ex:17:3}{$B=\log_2N+P\log_2P+(1-P)\log_2\frac{1-P}{N-1}=4.87$, $4.37$, $3.29$~біта на знак, тобто $7.3$, $6.6$, $4.9$~біт/с; для $10$~біт/с треба $2.3$~знака на секунду.}
\answer{ex:17:4}{Дисперсія похибки $2b/(Nm^2T)+\sigma_\xi^2$ на компоненту; внески рівні за $N\approx90$; границя $\tfrac12\log_2(1+0.25/0.04)\approx1.4$~біта на компоненту за вікно.}
\answer{ex:17:5}{Пара збігається за $\eta_bd^2+\eta_db^2<2$, тобто приблизно за $\eta<1$: за $0.8$ збігається, за $1.1$ стійкі коливання періоду $2$, за $1.5$ неперіодичні, за $2$ рознос; швидкості учнів треба розводити.}
\answer{ex:17:6}{Поріг $\approx4.4\sigma$; $102$~кбіт/с, $0.10$~мВт проти $205$~мВт для сирого потоку, у $2000$ разів менше; насичується (більше $3$ імпульсів) лише $5.7\cdot10^{-5}$ відліків.}
\answer{ex:17:7}{Відкрита задача. Близько $46$~біт/с за $\text{SNR}\approx0.9$ і близько $330$~біт/с за незалежного шуму. Якщо заважає шум, схожий на сигнал, витягнута інформація насичується зі зростанням кількості каналів; якщо незнання коду, вона зростає з кращим декодером (наприклад, що враховує часи імпульсів).}
